\documentclass[../main.tex]{subfiles}
\begin{document}
\section{Literature Review}

The literature on regression influence studies how fitted quantities depend on individual observations and groups of observations. Classical diagnostics characterise residual unusualness, leverage, and single-case influence; joint-deletion methods examine dependence on subsets; robust regression targets estimator stability under contamination. More recent data-dropping and influential-set methods formulate sensitivity directly as a subset-selection problem. Together, these strands motivate a coefficient-targeted analysis of how sensitivity is distributed across observations.

Outlyingness and target-coefficient influence provide the starting distinction. Outlyingness describes observations that are unusual relative to the sample, the data-generating process, or the fitted model. Target-coefficient influence describes the change in a prespecified regression coefficient induced by deleting an observation or set. Unusual observations can exert little influence on the target coefficient, while substantial coefficient influence can arise without conventional contamination outliers. The broader question is therefore how target-coefficient sensitivity becomes concentrated in particular subsets and how that structure depends on the underlying mechanism.

\subsection{Classical influence diagnostics}
\label{subsec:classical_influence_diagnostics}

Classical influence analysis distinguishes residual unusualness, leverage, and realised influence on an estimated regression quantity. Residual outliers are unusual relative to the fitted conditional mean, while leverage describes unusual covariate positions. Influence concerns the resulting change in an estimate or fitted quantity. These properties can occur separately, so an observation can be unusual with little target-coefficient influence or influential without an extreme residual \parencite{hampel1974influence,chatterjee1986influential}.

Case deletion provides the operational basis for regression influence diagnostics \parencite{cook2000detection}. Cook's distance summarises overall coefficient movement after deleting one observation, leverage describes the corresponding covariate geometry, and DFBETAS measures the standardised change in a selected coefficient \parencite{hoaglin1978hat,chatterjee1986influential}. DFBETAS is therefore closest to the coefficient-targeted perspective used in this thesis, although it remains a single-observation diagnostic.

Single-case deletion characterises sensitivity that is concentrated in individual observations. When several observations jointly support the same fitted relationship, deleting one member can leave most of that configuration intact. The resulting casewise coefficient change can remain moderate even when deletion of the group produces a substantial shift. Joint sensitivity therefore requires analysis at the level of observation sets.

\subsection{Joint influence, masking, and swamping}
\label{subsec:joint_influence}

Joint influence arises when the effect of a group cannot be inferred from separate deletion of its members. Masking occurs when jointly influential observations support one another strongly enough that individual deletion produces only modest diagnostic values \parencite{atkinson1986masking,rousseeuw1990unmasking}. Swamping describes the complementary case in which distortion of the fitted model causes observations outside the influential group to receive large diagnostic values \parencite{rousseeuw1990unmasking,hadi1993procedures}. Both phenomena make the subset, rather than the individual observation, the relevant diagnostic unit.

Early subset-deletion procedures addressed this problem through group deletion, sequential subset construction, robust distances, and influence-matrix methods \parencite{andrews1978finding,hadi1992identifying,hadi1993procedures,pena1995detection}. These methods established that regression influence can be distributed across observations. Their selected groups can nevertheless depend on screening rules, starting values, or sequential inclusion decisions.

The search problem becomes combinatorial once the deletion size exceeds one. For fixed \(k\), there are \(\binom{n}{k}\) possible deletion sets. Sequential and locally approximated procedures reduce this burden, but preliminary observation-level rankings can still inherit masking \parencite{pena2005new,zhao2019multiple,kuschnig2021hidden}. Modern influential-set methods therefore formulate joint deletion directly as optimisation over subsets of a prespecified size.

\subsection{Worst-case data dropping and most influential sets}
\label{subsec:worst_case_data_dropping}

Modern data-dropping analysis asks how an empirical conclusion changes when selected observations are removed. Influence-function approximations provide scalable estimates of these changes by expanding locally around the full-sample estimator, and extensions incorporate group effects and higher-order interactions \parencite{koh2017understanding,koh2019accuracy,basu2020second,fisher2023influence}. Related robustness analyses study the smallest deletion capable of changing a reported conclusion, such as a coefficient sign or significance statement \parencite{giordano2026automatic1,giordano2026automatic2}.

Searching directly over deletion sets is computationally difficult because a fixed-\(k\) audit contains \(\binom{n}{k}\) candidates. Practical methods therefore use greedy searches, relaxations, local approximations, and problem-specific optimisation \parencite{price2022hardness,suriyakumar2022algorithms,moitra2022provably,freund2023towards,hu2024most}. Approximation quality matters when the identities of the selected observations carry substantive meaning. Masking, interactions, curvature, and near-singular deleted designs can cause approximate rankings to differ from realised joint-deletion effects \parencite{huang2025approximations,rubinstein2025robustness}. This motivates coefficient-targeted optimisation when the objective is to characterise the most influential size-\(k\) subset.

\subsubsection{Fixed-\texorpdfstring{$k$}{k} most influential sets}
\label{subsubsec:fixed_k_most_influential_sets}

Most Influential Sets provide a coefficient-targeted formulation of joint deletion. \textcite{kuschnig2021hidden} show that groups of observations can remain concealed under conventional diagnostics while exerting substantial joint influence on an estimated coefficient. Fixed-\(k\) MIS searches jointly over deletion sets of a prespecified size, and subsequent work exploits linear-regression structure to obtain global optimisation of the stated fixed-residualised objective under appropriate conditions \parencite{konrad2026finding}. The resulting set has an objective-specific interpretation: among admissible sets of size \(k\), it maximises movement in the prespecified target coefficient under that formulation.

For a prespecified subset size \(k\), the fixed-\(k\) MIS objective can be represented conceptually as
\[
S_k^*
\in
\arg\max_{\substack{S\subseteq\{1,\ldots,n\}\\ |S|=k}}
\left|
\widehat{\beta}_{j,-S}
-
\widehat{\beta}_{j}
\right|.
\]
The parameter \(k\) specifies the deletion budget. Different values of \(k\) therefore define different sensitivity questions and can produce different influential sets. The fixed-residualisation convention, scope of global optimisation, and relationship to full-refit validation are defined in \autoref{sec:methodology}.

Optimisation and statistical calibration are separate tasks. A fixed-\(k\) maximum is selected from many candidate subsets, so its magnitude cannot generally be interpreted as though the deletion set had been prespecified. Where calibration is required in the simulations, the reference distribution is constructed separately from the optimisation step.

\subsection{Robust regression as a complementary objective}
\label{subsec:robust_regression}

Robust regression addresses coefficient stability by changing the estimation rule used in the presence of contamination or other departures from the reference model. M-estimation limits residual influence through alternative loss functions, while high-breakdown procedures such as Least Trimmed Squares and MM regression use resistant fitting criteria to maintain stability under more substantial contamination \parencite{huber1964location,huber1973robust,rousseeuw1984applying,yohai1987high}. These procedures are therefore designed primarily for stable estimation.

Coefficient-targeted influence analysis and robust regression characterise different dimensions of sensitivity. Fixed-\(k\) MIS identifies the subset whose joint deletion produces the largest movement in a prespecified coefficient, while robust estimators seek coefficient stability under contamination. Their performance criteria therefore need not coincide. This difference motivates the localisation and estimation comparison in \autoref{subsec:estimation-performance}.

\subsection{Model misspecification and coefficient sensitivity}
\label{subsec:model_misspecification}

Model misspecification provides a potential source of target-coefficient sensitivity without requiring a conventional contamination set. Under conditional-mean misspecification, least squares can converge to a well-defined linear approximation even when the fitted regression differs from the population regression function \parencite{white1980using}. The finite-sample influence of particular observations then depends on how the departure interacts with the covariate distribution and the target coefficient. Misspecification can therefore alter both the coefficient itself and the way its sensitivity is distributed across observations.

Different departures can generate different patterns of coefficient sensitivity. Omitted variables, nonlinear functional form, missing interactions, heterogeneous slopes, and threshold behaviour alter the conditional-mean relationship, while endogeneity changes the relation between regressors and disturbances and heteroskedasticity changes the conditional variance. Specification-focused procedures such as RESET and the Breusch--Pagan test target particular implications of these departures \parencite{ramsey1969tests,breusch1979simple}. Coefficient-targeted deletion instead asks whether the departure is expressed through concentrated influence on the prespecified coefficient.

The misspecification simulation therefore uses MIS to study how model departures are expressed through finite-sample target-coefficient sensitivity. The analysis asks whether different forms of misspecification generate concentrated joint influence, whether identifiable affected subgroups coincide with the selected influential set, and where the sensitivity response becomes weak or invariant. Specification testing and model correction address separate properties of the fitted model.

Together, these literatures motivate a study of the structure of joint target-coefficient sensitivity. Classical diagnostics establish the distinction between unusualness and realised influence, subset-deletion methods make groups of observations the relevant diagnostic unit, robust regression separates localisation from estimator stability, and misspecification broadens the mechanisms through which concentrated sensitivity can arise. The methodology below operationalises this perspective through fixed-\(k\) coefficient-targeted influential sets.
\end{document}